\documentclass[10pt,reqno]{amsart}

\usepackage[T1]{fontenc}
\usepackage[utf8]{inputenc}
\usepackage{newtxtext,newtxmath}
\usepackage{microtype}
\usepackage{mathtools}
\usepackage{booktabs}
\usepackage{array}
\usepackage{enumitem}
\usepackage{xcolor}
\usepackage{hyperref}
\usepackage[a4paper,margin=28mm,headheight=15pt]{geometry}

\definecolor{LinkBlue}{HTML}{174A7E}
\hypersetup{
  colorlinks=true,
  linkcolor=LinkBlue,
  citecolor=LinkBlue,
  urlcolor=LinkBlue
}

\setlength{\parskip}{2pt}
\setlist[itemize]{leftmargin=6mm,itemsep=1.5pt,topsep=3pt}
\setlist[enumerate]{leftmargin=7mm,itemsep=2pt,topsep=3pt}
\renewcommand{\arraystretch}{1.15}
\emergencystretch=2em

\newtheorem{theorem}{Theorem}[section]
\newtheorem{proposition}[theorem]{Proposition}
\newtheorem{lemma}[theorem]{Lemma}
\newtheorem{corollary}[theorem]{Corollary}
\theoremstyle{definition}
\newtheorem{definition}[theorem]{Definition}
\newtheorem{remark}[theorem]{Remark}

\title[An Explicit Four-Point Rational Counterexample to AIM Question 47]
{An Explicit Four-Point Rational Design Counterexample to AIM Question 47 in Algebraic Statistics}
\author{Argus Open Research Project}
\date{August 27, 2026}
\subjclass[2020]{Primary 62R01; Secondary 13P10, 14Q99}
\begin{document}
\begin{abstract}
Question~47 of the AIM problem list in computational algebraic statistics asks whether reducing the minimal generators $m_1,\dots,m_r$ of a monomial ideal $M\subset\mathbb{Q}[x_1,\dots,x_n]$ modulo a finite design $D=\{v_1,\dots,v_N\}\subset\mathbb{Q}^n$ against its standard monomial basis $\mathcal{O}=\mathrm{Std}(M)$ yields polynomials $f_1,\dots,f_r\in I(D)$ whose common zero locus $V(f_1,\dots,f_r)$ equals $D$. We give an explicit four-point rational counterexample in $\mathbb{Q}[x,y]$ where $V(f_1,f_2)$ contains the entire line $x=y$.
\end{abstract}
\maketitle
\section{The Four-Point Counterexample}
In $\mathbb{Q}[x,y]$, define the four-point rational design and monomial ideal
\[
  D \;=\; \bigl\{(0,0),\ (1,1),\ (2,2),\ (0,1)\bigr\} \subset \mathbb{Q}^2,
  \qquad
  M \;=\; \langle x^2, y^2\rangle,
  \qquad
  \mathcal{O} \;=\; \mathrm{Std}(M) \;=\; \{1, x, y, xy\}.
\]
\begin{theorem}
For $(D, M, \mathcal{O})$ above:
\begin{enumerate}[label=\textup{(\arabic*)}]
  \item Evaluating $(1,x,y,xy)$ on the four points of $D$ yields the $4\times 4$ matrix with determinant $-2\neq 0$, so $\mathcal{O}$ is a basis of $\mathbb{Q}[x,y]/I(D)$.
  \item The unique interpolants of the minimal generators $m_1=x^2$ and $m_2=y^2$ in $\operatorname{span}_{\mathbb{Q}}\mathcal{O}$ on $D$ are $x^2\equiv xy$ and $y^2\equiv xy-x+y\pmod{I(D)}$, producing
  \[
    f_1 \;=\; x^2 - xy \;=\; x(x-y),
    \qquad
    f_2 \;=\; y^2 - xy + x - y \;=\; (y-1)(y-x).
  \]
  \item The common zero locus is $V(f_1,f_2) = \{(t,t) : t\in\mathbb{Q}\}\cup\{(0,1)\}\supsetneq D$; in particular $(3,3)\in V(f_1,f_2)\setminus D$.
\end{enumerate}
\end{theorem}

\section{Methodological Framework and Verification Invariants}

The mathematical analysis for Result~20001224 in the domain of algebraic statistics and computational commutative algebra follows a strict separation between the exact statement decided and nearby formulations with stronger hypotheses. Specifically:
\begin{enumerate}[label=\textup{(\alph*)}]
  \item \emph{Statement alignment:} Every quantifier, ambient category, normalization convention, and regularity or finiteness hypothesis is fixed against the primary problem source prior to evaluating candidate constructions or obstructions.
  \item \emph{Exactness of algebraic and combinatorial steps:} All finite-dimensional linear-algebraic, polynomial, lattice, or graph-theoretic invariants appearing in the proof above are established in exact rational, integral, or symbolic arithmetic without floating-point approximation.
  \item \emph{Classification:} Within the program taxonomy, this result is categorized as a \emph{confirmed counterexample} deciding the scope: D=\{(0,0),(1,1),(2,2),(0,1)\} in Q\textasciicircum\{\}2 with M=<x\textasciicircum\{\}2,y\textasciicircum\{\}2> yields V(f\_1,f\_2) = V(x-y) U \{(0,1)\} != D.
\end{enumerate}

\section{Scope and Open Boundaries}

To prevent over-claiming beyond the proved theorem, we record the explicit boundaries that are \emph{not} decided by the argument above:
\begin{itemize}
  \item full border basis interpolation.
\end{itemize}

\begin{thebibliography}{9}
\bibitem{PrimarySource}
Primary problem list and source literature for Result~20001224 (algebraic statistics and computational commutative algebra), archived with exact theorem locators and verification certificates in the accompanying repository package.
\end{thebibliography}

\end{document}
